\section{Introduction}
\label{sec:introduction}

Storage-backed inference repeatedly brings selected weights into an execution
cache rather than keeping every execution weight resident. This approach can
reduce device residency, but a small cache alone does not bound the resources
used to fill it. A read may retain host buffers, an upload may outlive its
caller, and an evicted binding may still have device consumers. Releasing a
logical request therefore does not necessarily make its capacity reusable.
The central problem is to coordinate weight materialization, resource admission,
and ownership through the last consumer.

Weight streaming and prefetch are established techniques in systems such as
FlexGen, ZeRO-Inference, and PowerInfer-2~\cite{sheng2023flexgen,deepspeed2022zeroinference,xue2024powerinfer2}.
Expert caches and speculative loading also address sparse demand in MoE
models~\cite{xue2024moeinfinity,eliseev2023fastmoeoffloading}.
SGLang also provides a specialized storage-backed expert path in the inspected
source~\cite{sglang2026expertpack}.
\system{} addresses a complementary architectural question: how should a
materialization schedule bind source identity, stage-local accounting, and
completion evidence so that capacity is not reused prematurely? We distinguish
three properties that are often conflated: a configured resource budget, a
measured physical working set, and evidence that a particular operation has
finished using its resources. Neither storage offload itself nor performance
leadership is the claimed novelty.

The design follows \emph{plan, admit, materialize, publish, execute, retire}.
A frozen plan identifies the source and destination generations. Admission
acquires the claims required by each stage. Completion validates the exact
operation before publication, while the owner retains buffers and charges
until the relevant consumers finish. If completion cannot be proved, the
intended discipline retains uncertain resources instead of treating an error
or cancellation as permission to reuse them. The resulting conservation model
is conditional on complete ownership and transfer rules; the serving evidence
observes selected host and expert-cache accounts, not every class in an
integrated system-wide ledger.

The implemented expert path separates read scheduling from transfer policy.
A persistent worker serves bounded checkpoint-read slots. The blocking schedule
waits at demand; lookahead submits one next-expert read before current expert
computation. Independently, grouped synchronous transfer batches the expert's
weight and scale copies behind one upload completion boundary. Both choices
preserve the expert order, arithmetic, and synchronous forward interface.
Worker lifetime, payload lifetime, upload completion, and device-consumer
retirement remain distinct: an idle live worker can still own charged metadata
and stack capacity.

\paragraph{Contributions.}
\begin{enumerate}
  \item An architecture for bounded storage-backed weight residency that connects
  source identity, stage-local admission, publication, and retirement, while
  distinguishing logical accounting from physical memory and producer completion
  from consumer quiescence.
  \item Reference algorithms for exact materialization and bounded cleanup,
  together with an implemented persistent read worker, bounded next-expert
  scheduling, and grouped synchronous expert uploads.
  \item A profile-scoped evaluation linking numerical and HTTP output checks,
  logical work and resource accounts, client timing, selected-process profiling,
  and component-failure and cancellation behavior, including negative results.
\end{enumerate}

\paragraph{Research questions.}
\begin{description}
  \item[RQ1: Bounded completion.] Can requests complete under the declared
  accounting policies, and how does residency change across coupled cache budgets?
  \item[RQ2: I/O demand.] What work reaches the logical reader and shared
  filesystem client, and which physical traffic remains unknown?
  \item[RQ3: Scheduling and transfer.] How do read scheduling and upload grouping
  affect observed waiting, synchronization, and client outcomes?
  \item[RQ4: Failure and custody.] What resources and state remain after
  clean abort, uncertain completion, and HTTP cancellation?
\end{description}

The core serving workload uses Qwen3.5-35B-A3B-FP8 through the public CLI,
with BF16 execution weights and explicitly defined F32 boundaries, not a
full-model W8A16 quality claim. Four acceptance servers complete 48 checked
requests. A separate grouped-versus-synchronous transfer pilot completes eight
fresh server lifecycles and 96 requests, including 40 measured eight-token
quality checks. With lookahead, grouped uploads reduce observed client latency
and increase phase-qualified goodput; with blocking reads, time to first content
regresses despite nearly unchanged end-to-end latency. There are only two
pairs per read schedule. An earlier five-pair read-scheduling study instead
compares blocking with lookahead under its own fixed implementation, finding
less demand waiting and lower observed client latency with equal logical reads.
These are separate comparisons, not a pooled population or evidence of a
universal improvement.

The evaluation establishes neither a physical capacity frontier nor a matched
external-system ranking. DeepSeek-V4 end-to-end execution is untested;
architecture-level applicability does not qualify another model adapter.
Table~\ref{tab:evidence} consolidates the measurement boundaries.
Section~\ref{sec:scope} defines the model and assumptions.
Section~\ref{sec:design} develops the mechanisms and algorithms, and
Section~\ref{sec:implementation} maps them to the implemented path.
Section~\ref{sec:eval-setup} fixes the methodology before
Section~\ref{sec:evaluation} answers the four questions. Related work,
limitations, and separate historical evidence follow.
